% AUTO-GENERATED durch diagram_generator (P3, Latenz-Verteilung: p50--p99-Spanne, Punkt+Whisker)
% EHRLICH: das WIDE-Schema traegt NUR p50/p99 je Permutation, NICHT die rohen Einzel-Op-Latenzen.
% Punkt = nearest-rank-Median der p50-Werte; plus-Whisker hoch bis Median der p99-Werte. KEIN Box-Plot
% (nur 2 Perzentile -> Quartile waeren erfunden = Phantom). 1 addplot je op-Art (Farbe+Legende).
\begin{figure}[!htbp]
\centering
\resizebox{\textwidth}{!}{%
\begin{tikzpicture}
\definecolor{rangeop0}{RGB}{217,82,82}
\begin{axis}[
    width=0.7800\textwidth,
    height=0.4000\textheight,
    scale only axis,
    title={Latenz-Spanne p50-p99 je Suchalgorithmus x Operation},
    xlabel={Suchalgorithmus x Operation},
    ylabel={Latenz (ns, log; Punkt=p50, Whisker bis p99)},
    grid=both,
    enlargelimits=0.05,
    ymode=log,
    symbolic x coords={k\_ary/lookup},
    xtick=data,
    x tick label style={rotate=60,anchor=east,font=\tiny},
    legend pos=north west,
    legend style={font=\tiny,legend cell align=left},
    mark size=2.4pt,
]
\addplot[only marks,mark=*,color=rangeop0,
    error bars/.cd, y dir=plus, y explicit, error bar style={line width=0.7pt,color=rangeop0}]
coordinates {
    (k\_ary/lookup,750.0000) +- (0,420.0000)
};
\addlegendentry{lookup}
\end{axis}
\end{tikzpicture}
}%
\caption{Latenz-Spanne p50-p99 je Suchalgorithmus x Operation; Punkt=p50, Whisker=p99.}
\end{figure}
